\section{Results, input models, and the remaining objective}
\label{sec:scope}

The most useful next step for this project is to change which intermediate
objects are materialized. The existing recognizer already avoids much
unnecessary work through local simplification, algebraic obstructions,
connected-sum structure, and improved Khovanov computations. Those changes
do not make an explicit complex small merely because its boundary is small.
The present contribution supplies an exact arithmetic route around that
problem on a well-defined family, and a local certificate route that applies
to a larger family of diagrams.

The central engineering recommendation is to retain certified rational
pieces as primitive integer slopes and decide supported gluings arithmetically.
When a general diagram contains a tangle that is incompatible with every
unknotting complement, a checked local occurrence can end the decision
without evaluating the exterior. Each recommendation has a precise input
contract, a proof obligation, executable code, and a measured limitation.

\subsection{What is established here}

Table~\ref{tab:claims} distinguishes mathematical conclusions from
implementation and discovery claims. This distinction is necessary because
a fast algorithm on a supplied presentation is not automatically an algorithm
that discovers that presentation in an arbitrary diagram.

\begin{table}[htbp]
\centering\small
\begin{tabularx}{\textwidth}{L{0.25\textwidth}X L{0.21\textwidth}}
\toprule
Contribution & Exact scope & Status\\
\midrule
Exponential continuation profiles &
Four marked endpoints, \(O(m)\) crossings per prefix and continuation,
all cross closures knots &
Proved; arithmetic diagnostic through \(m=10\)\\
\addlinespace
Exponential explicit prefix complex &
Any full delooped Bar--Natan complex for a specified rational prefix
inside an unknot family &
Proved; not a scan-order-independent decision lower bound\\
\addlinespace
Montesinos arithmetic classifier &
Supplied finite-slope continued fractions and integer term; one-component
input checked &
Implemented, complete on this input class\\
\addlinespace
Forbidden subtangle search &
Literal connected disk patterns in a validated plane knot diagram;
four distinct cut edges &
Implemented for supplied or fixed finite catalogue\\
\addlinespace
Quasi-polynomial discovery &
Literal admissible patterns with at most \(O(\log n)\) continued-fraction
coefficients &
Enumeration algorithm proved; exhaustive enumerator not implemented\\
\addlinespace
General quasi-polynomial recognizer &
Every valid \(n\)-crossing knot diagram &
Not established by this contribution\\
\bottomrule
\end{tabularx}
\caption{Claim ledger. \emph{Proved} refers to the mathematical arguments in
this report, rather than formal verification in a proof assistant.}
\label{tab:claims}
\end{table}

\subsection{Three input sizes that must be kept separate}
\label{subsec:input-sizes}

Let \(n\) be the number of crossings of a supplied planar diagram (PD).
Its four-valent plane map, crossing information, and component structure
are part of the input. With canonically bounded labels its encoding has
\(O(n\log(n+2))\) bits. If users supply labels with arbitrarily many bits,
their actual encoding length must also be charged to parsing. Word-operation
bounds below assume \(O(\log(n+2))\)-bit indices after relabeling.

Let \(B\) be the binary length of a structured tangle source. It includes
the signed integer term, all signed continued-fraction coefficients,
and list delimiters. Counting delimiters matters: a list containing many
zeros can have almost no twists but still take time to read and evaluate.
Let
\[
 C=|e|+\sum_{i,j}|a_{i,j}|
\]
be the number of crossings in the literal expansion of that source.
The compiler in this contribution constructs precisely that literal diagram,
before any optional Reidemeister simplification. In general \(C\) can be
exponential in \(B\), so an \(O(B^2)\) arithmetic result does not justify
expanding a PD in \(O(B^2)\) time.

Finally, let \(h\) denote the \emph{number} of continued-fraction coefficients
in a local template description. The discovery theorem in
Section~\ref{sec:mon:montesinos} assumes both a bound on \(h\) and an explicit
crossing budget. It never equates coefficient count with bit length.

We use \emph{quasi-polynomial} for a bound
\(\exp((\log n)^{O(1)})\). The more specific target
\(n^{O(\log n)}=\exp(O((\log n)^2))\) is stronger than
\(n^{O((\log n)^2)}=\exp(O((\log n)^3))\).
Whenever an exponent matters we write it explicitly.

\subsection{The audited repository baseline}

The work starts from ProveIt commit
\code{86f23490006c880cd01db9bec10761d39c6395c8}, retrieved on 8 October
2026~\cite{proveit-baseline}. The audited executable path is the
\code{fast} subdirectory under \code{Topology/UnknotRecognition}. Its exact source and relevant tests,
examples, and recent research documents were retrieved through the GitHub
interface. Git blob hashes were checked against the fetched bytes for
82 files; the source manifest accompanies the integration package.
The baseline test suite contains 134 tests, despite the older test count
still present in an earlier README passage.

The baseline already includes Reidemeister reduction, descending-diagram
recognition, complete specialized three-braid decisions, signed Seifert
and Rasmussen-style obstructions, modular and exact Alexander paths,
a Jones filter, visible connected-sum and interlacement decomposition,
and several exact Khovanov backends. Component sharing, saturated
multiplicities, Euler bounds, sparse maps, and composition caching are
existing work. The recent reports additionally discuss radical
minimization, exact homological windows, continuation shortening, Fitting
splits, terminal boundary tests, twist macros, and hierarchy restart counts.
These are context for the present work and are not relabeled as new
contributions.

The new production stage is deliberately earlier than complex construction
when a checked rational source exists. Its mathematical information is
different from an approximate Jones or Alexander value: within the supported
Montesinos class it is a complete decision. The general exact fallback
continues to have an exponential worst-case upper bound and a finite-budget
\unknown\ result. The JSON flag reporting the absence of a general
quasi-polynomial guarantee is retained.

\subsection{Relevant literature and the status of the broad target}
\label{subsec:literature-status}

The exact Khovanov decision criterion ultimately uses the theorem of
Kronheimer and Mrowka that reduced Khovanov homology detects the
unknot~\cite{kronheimer-mrowka}. Over \(\F_2\), a reduced rank of one also
forces rational reduced rank one by the universal coefficient inequality.
This supplies a sound decision criterion; it does not give a fast method
for computing the required rank.

Lackenby's March 2021 Oxford slides announce a general unknot algorithm
with running time \(2^{O((\log n)^3)}\)~\cite{lackenby-2021}.
The power three in the logarithm should not be silently replaced by two.
The July 2026 manuscript \emph{Incompressible surfaces, hierarchies and
unknot recognition} gives a new hierarchy algorithm and substantial
topological structure~\cite{lackenby-2026}. Section 9 of that manuscript
explicitly separates an iteration bound from a full running-time analysis
because several operations have not been specified sufficiently for the
latter. Its displayed iteration estimate is
\[
 L(g+1)^L,
\]
where \(L\) is the maximum hierarchy length and \(g\) bounds surface
pattern-complexity. The refined construction in Section 10 controls
hierarchy length in terms of the initial handle complexity.

This report uses that distinction constructively. A project implementation
needs explicit operation algorithms, encoding-size bounds, and the
appropriate logarithmic hierarchy estimate before inheriting a
quasi-polynomial bound. A 2021 announcement, a 2026 iteration estimate,
and a proved executable bound for the present code are different kinds
of evidence. We do not infer the last from either of the first two.
This is a statement about the sources and implementation audited here,
not a claim to have ruled out every other possible algorithm in the
literature.

The rational-tangle classification and gluing formulas used below are
classical; we use the treatments of Kauffman and
Lambropoulou~\cite{kauffman-lambropoulou-tangles,
kauffman-lambropoulou-knots,KauffmanLambropoulouHard}.
The local obstruction is the theorem of Nogueira and
Salgueiro~\cite{nogueira-salgueiro}. The lower bounds in the next sections
are project-specific deductions from these tools and the functorial
Khovanov framework, with their own proofs.

\subsection{Use of the suggested mathematics repository}

The suggested \code{openai/math} repository was inspected at the following commit:
\begin{center}
\code{adc7f1241b42e322a6451854ab7e4b4c146bf78a}.
\end{center}
Its description distinguishes
research manuscripts at different stages from formally verified results.
The relevant inspiration was a standard context-separation proof pattern:
equal summaries must give equal answers after every allowed continuation.
The inspected section applies this idea to word relations and
automata~\cite{OpenAIMathContextSeparation}.

Our rational-tangle construction below is independent of that manuscript's
headline automata claim. It specifies actual knot diagrams, proves the
diagonal acceptance matrix using unimodular arithmetic and disjoint cones,
and gives a separate arithmetic representation. No theorem from the
suggested repository is used as a black box for topology or complexity.
The exact source path, commit, and blob hash are recorded in the package's
research provenance.
